\documentclass{article}
\usepackage[dblblindworkshop]{neurips_2026}
\workshoptitle{Muslims in ML}
\usepackage[utf8]{inputenc}
\usepackage[T1]{fontenc}
\usepackage{hyperref}
\usepackage{url}
\usepackage{booktabs}
\usepackage{amsmath,amsfonts}
\usepackage{microtype}
\newcommand{\taskA}{Task~A}
\newcommand{\taskB}{Task~B}
\title{Checking Quran Recitation from Noisy ASR Transcripts:\\A Competition with Agent-Assisted Annotation}
\author{Anonymous Author(s)}
\begin{document}
\maketitle
\begin{abstract}
We propose two tasks for checking Quran recitation from ASR: detecting and splitting ayahs, and labeling localized transcript/reference events. Task A has a public 258-recording benchmark. For Task B we completed a human-approved annotation set of 100 recording cases, 314 ayah chunks and 115 within-ayah events, froze an executable scorecard over combined labels and spans, and evaluated three baselines. They locate differences accurately and then misjudge them: a plain diff calls 88.5\% of benign and corrected events mistakes, the deployed pipeline 53.8\%. A 30-record sample accompanies this submission. An agent track will annotate a separate unlabelled pool and build an algorithm for private scoring. We document the input exposure that prevents claiming an unseen-input final test.
\end{abstract}

\section{Task and motivation}
Quran memorization involves a listener checking recitation. An ASR-based application must handle opening formulas, repeats, restarts, repairs and missing or substituted words, which edit distance alone does not distinguish. Our labels assess transcript/reference differences, not whether speaker or ASR caused them; unwritten vowel and tajweed errors are out of scope.

In \textbf{\taskA\ (detect and split)} a full transcript maps to ayah IDs and word assignments, or abstains on non-Quran input; the public v1.1 benchmark exists. In \textbf{\taskB\ (event annotation)} outputs carry combined labels such as \texttt{substitution\_mistake} and \texttt{repetition\_benign} with hypothesis and reference spans. Review covers whole recordings and opening formulas, and algorithms receive reviewed chunks, ayah IDs and exact references, so scoring excludes detection errors.

\section{Data and findings from human review}
The public \taskA\ release has 258 transcripts: 254 Quranic recordings, four non-Quran negatives and 1{,}267 reviewed chunks covering 604 ayahs from 45 surahs, split 151/107 recordings and 711/556 chunks, duplicates kept in one partition.

\textbf{Completed \taskB\ review.} All \textbf{100 selected recording cases} are human-approved under rubric v0.19: 314 chunks, 233 no-event, 115 within-ayah events and 47 opening-formula annotations, over 274 ayahs in 38 surahs and ten learner IDs, one contributing 46 recordings; three approved reserves sit outside the scored set. These purposive counts do not estimate prevalence or imply generalization to unseen reciters. A deterministic 30-record sample covering every label accompanies this submission.

We verified recording IDs and transcripts against both source exports and excluded punctuation-normalized duplicates and redundant clean subpassages.

Four rules settled during review. \textbf{Omissions} include missing beginnings and endings, with no exemption for a final ayah. \textbf{Repeats} span the original phrase and its extra copy against one reference copy; a wrong-then-correct attempt is corrected, a correct-then-incorrect restatement is \texttt{substitution\_mistake} spanning both. \textbf{Spelling} differences removed by agreed normalization receive no event; residual forms use \texttt{spelling\_benign}, which permits no blanket final-letter deletion. \textbf{Localization} uses verbatim words and half-open spans in the original token arrays, with empty hypothesis spans for omissions. Opening formulas and letter names are benign, and possible ASR origin does not create an uncertain verdict.

\section{Metric and baselines}
The \taskA\ scalar sums detection, split and abstention errors; earlier work reported held-out 0.760 for the incumbent pipeline and 0.079 for the best of six coding-agent runs.\footnote{\emph{Autoresearch with coding agents: generalizers and metric-maximizers on Quran recitation data}, AIST 2026.}

\textbf{\taskB\ scorecard.} Predicted and gold events are matched one-to-one within a chunk by maximum span similarity, so one broad prediction cannot claim several gold events. Similarity averages hypothesis- and reference-span overlap, with empty spans as omission anchors and insertion points matched by proximity. With $\text{miss}$ the share of gold mistakes no mistake prediction matched, $\text{false\_alarm}$ predicted mistakes matching no gold event, $\text{benign}$ gold benign or corrected events called mistakes, and $\text{clean}$ no-event chunks carrying any flag, $\texttt{study3\_v19\_score} = 2\cdot\text{miss} + \text{false\_alarm} + \text{benign} + \text{clean}$: a missed mistake fails the learner and is weighted double, an extra flag costs a moment of review. Label agreement and span IoU over matched pairs are reported, not summed. Predicting nothing scores 2 and the gold annotation 0; oracle tests confirm both endpoints and that one chunk-wide flag leaves most gold mistakes in multi-event chunks unmatched. The evaluator is frozen before any agent run.

\begin{table}[h]\vspace{-10pt}\centering\footnotesize\setlength{\tabcolsep}{4.5pt}
\begin{tabular}{lrrrrrr}\toprule
Baseline & score & miss & false alarm & benign & clean & span IoU\\\midrule
Predict nothing & 2.000 & 1.000 & 0.000 & 0.000 & 0.000 & --- \\
Naive diff & 0.988 & 0.045 & 0.009 & 0.885 & 0.004 & 0.921 \\
Incumbent pipeline & 0.768 & 0.079 & 0.068 & 0.538 & 0.004 & 0.883 \\
\bottomrule
\end{tabular}\vspace{-5pt}
\caption{\taskB\ baselines on 314 reviewed chunks: 233 no-event, 89 mistakes, 26 benign or corrected.}\vspace{-12pt}
\end{table}

Neither baseline has trouble locating differences: they miss 4--8\% of mistake events, localize matched events at 0.88--0.92 IoU and almost never disturb a clean chunk. They fail on the distinction the task exists for. A plain diff calls 88.5\% of benign and corrected events mistakes, and the deployed pipeline, which already models repeats and restarts, still calls 53.8\%, with localization recall 1/5 on repetitions. The headroom is in judgment, not detection.

The legacy \taskB\ pilot of 99 machine-labeled chunks under the old schema does not validate the new gold; its documented export mismatch is preserved in the historical release. Budgets and comparison conditions still need freezing.

\section{Competition plan and exposure controls}
The agent track supplies unlabelled development transcripts, the Quran reference, definitions and disjoint examples; entrants produce annotations and an algorithm, and the final code is scored, so gold performance assesses the combined process rather than annotation accuracy alone. The development bundle excludes selected inputs and duplicate-equivalent material, and a clarification exercise precedes measured runs. Agents receive no gold answers and no score feedback; reviewer files, source repositories, answer-bearing history and annotation-session memory stay out of agent runtimes, and shared documentation carries aggregate findings and synthetic examples only.

\textbf{Exposure limitation.} All 100 selected transcripts match public \taskA\ inputs and twelve appear in the legacy pilot: the adjudications and selection membership are private, the inputs are not. Runtime isolation cannot establish absence of prior exposure, so an unseen-input final ranking needs unreleased recordings; a read-only production check has since found 63, of which 57 are textually distinct from the released export.\end{document}
